% ===========================================================================
% Figure 1 -- motivating example
%
% Two separate background panels: (a) document, (b) code.
% Three coloured frames, one per inconsistency:
%   1 blue   -- both ends in the document  (inter-document layer)
%   2 orange -- document to code           (document-code layer)
%   3 teal   -- document to code           (document-code layer)
%
% Geometry is driven by explicit coordinates so left/right panels and all
% three frames are perfectly aligned rectangles.
% ===========================================================================
\begin{figure}[t]
  \centering
  \begin{tikzpicture}[
      font=\footnotesize,
      tt/.style={font=\ttfamily\scriptsize, align=left, text=black,
                 inner sep=0pt, outer sep=0pt},
      head/.style={font=\footnotesize\bfseries, text=black,
                   inner sep=0pt, outer sep=0pt},
      lab/.style={font=\scriptsize\bfseries, inner sep=0pt, outer sep=0pt},
    ]

  % -----------------------------------------------------------------------
  % Geometry constants (all in cm)
  % -----------------------------------------------------------------------
  % panels
  \def\PadO{0.22}      % outer padding around each panel
  \def\PadI{0.20}      % inner padding (content inset from panel edge)
  \def\Gap{0.55}       % gap between the two panels
  \def\Lw{6.8}         % left panel content width
  \def\Rw{5.6}         % right panel content width
  % derived
  \pgfmathsetmacro{\LpanL}{-\PadO}
  \pgfmathsetmacro{\LpanR}{\Lw+\PadO}
  \pgfmathsetmacro{\RpanL}{\Lw+\Gap+\PadO*2+\Gap*0}
  \pgfmathsetmacro{\RpanR}{\RpanL+\Rw+\PadO*2}
  % content x-anchors
  \def\Lx{0}
  \pgfmathsetmacro{\Rx}{\Lw+\Gap+\PadO+\PadI+0.18}

  % vertical anchors (top of each group label)
  \def\Ya{ 0.00}
  \def\Yb{-2.60}
  \def\Yc{-5.45}
  % bottom of last content (set after node placement via coordinate)

  % -----------------------------------------------------------------------
  % Headers
  % -----------------------------------------------------------------------
  \node[head, anchor=south west] at (\Lx, 0.42)
    {(a)\ \texttt{SKILL.md}\,---\,the document};
  \node[head, anchor=south west] at (\Rx, 0.42)
    {(b)\ \texttt{scripts/}\,---\,the code};

  % -----------------------------------------------------------------------
  % GROUP 1  inter-document layer  (blue, left panel only)
  % -----------------------------------------------------------------------
  \node[lab, text=blue!65!black, anchor=north west] (T1) at (\Lx, \Ya)
    {\textcircled{1}\enspace inter-document layer};
  \node[tt, text width=\Lw cm, anchor=north west] (L1)
    at ([yshift=-4pt]T1.south west){%
      :3\enspace "data visualizations with D3.js \colorbox{green!15}{(v6)}."\\[1pt]
      :75\ \ "Pin D3 to a specific version
              (default: \colorbox{red!15}{d3@7.9.0})."\\[1pt]
      \textcolor{black!45}{// one file, two contradictory major versions}};

  % -----------------------------------------------------------------------
  % GROUP 2  document-code layer  (orange)
  % -----------------------------------------------------------------------
  \node[lab, text=orange!80!black, anchor=north west] (T2) at (\Lx, \Yb)
    {\textcircled{2}\enspace document--code layer};
  \node[tt, text width=\Lw cm, anchor=north west] (L2)
    at ([yshift=-4pt]T2.south west){%
      :59\ "Sort input rows deterministically before binding to marks."\\[1pt]
      :60\ "\ldots\,stable grouping order
            (\colorbox{green!15}{Array.from(grouped.keys()).sort()})"\\[1pt]
      \textcolor{black!45}{// the row order must be stable and explicit}};
  \node[tt, text width=\Rw cm, anchor=north west] (R2)
    at ([yshift=-4pt]T2.south west -| \Rx pt,0){%
      bubble\_chart\_example.js:62\\[1pt]
      \colorbox{red!15}{const sectors =}\\
      \colorbox{red!15}{[...new Set(data.map(d => d.sector))]}\\[1pt]
      \textcolor{black!45}{// unsorted: order follows input}};

  % -----------------------------------------------------------------------
  % GROUP 3  document-code layer  (teal)
  % -----------------------------------------------------------------------
  \node[lab, text=teal!65!black, anchor=north west] (T3) at (\Lx, \Yc)
    {\textcircled{3}\enspace document--code layer};
  \node[tt, text width=\Lw cm, anchor=north west] (L3)
    at ([yshift=-4pt]T3.south west){%
      :68--71\ "if a force layout is required: fix the
                \colorbox{green!15}{tick count},
                fix initial positions \colorbox{green!15}{deterministically},
                run \colorbox{green!15}{exactly N ticks} and stop."\\[1pt]
      \textcolor{black!45}{// layout output must be reproducible}};
  \node[tt, text width=\Rw cm, anchor=north west] (R3)
    at ([yshift=-4pt]T3.south west -| \Rx pt,0){%
      bubble\_chart\_example.js:56\\[1pt]
      \colorbox{red!15}{const simulation =}\\
      \colorbox{red!15}{d3.forceSimulation(data)}\\[1pt]
      \textcolor{black!45}{// default tick count, no seed,}\\
      \textcolor{black!45}{// no fixed initial positions}};

  % -----------------------------------------------------------------------
  % Coloured inconsistency frames
  % Each frame spans from its label node to the bottom of its content,
  % left frame: Lx .. LpanR,  cross frame: Lx .. RpanR
  % We derive the bottom from the content nodes.
  % -----------------------------------------------------------------------
  % F1: left panel only
  \draw[blue!60!black, rounded corners=3pt, line width=0.8pt]
    ([xshift=-5pt, yshift=5pt]T1.north west)
    rectangle
    ([xshift=\LpanR pt-\Lx pt+5pt, yshift=-5pt]L1.south west
      |- [yshift=-5pt]L1.south east);

  % F2: full width
  \draw[orange!75!black, rounded corners=3pt, line width=0.8pt]
    ([xshift=-5pt, yshift=5pt]T2.north west)
    rectangle
    ([yshift=-5pt]R2.south east |- [yshift=-5pt]L2.south east);

  % F3: full width
  \draw[teal!60!black, rounded corners=3pt, line width=0.8pt]
    ([xshift=-5pt, yshift=5pt]T3.north west)
    rectangle
    ([yshift=-5pt]R3.south east |- [yshift=-5pt]L3.south east);

  % -----------------------------------------------------------------------
  % Background panels  (drawn last so they go behind frames)
  % We use the bounding boxes of the content nodes plus fixed padding.
  % -----------------------------------------------------------------------
  % Determine overall vertical extent from node positions
  \coordinate (top_y)  at ([yshift=+0.55cm]T1.north);
  \coordinate (bot_y)  at ([yshift=-0.28cm]L3.south);

  \begin{scope}[on background layer]
    % Left panel
    \fill[black!3, rounded corners=4pt]
      (top_y -| \LpanL pt,0) rectangle (bot_y -| \LpanR pt,0);
    \draw[black!22, rounded corners=4pt]
      (top_y -| \LpanL pt,0) rectangle (bot_y -| \LpanR pt,0);
    % Right panel
    \fill[black!3, rounded corners=4pt]
      (top_y -| \RpanL pt,0) rectangle (bot_y -| \RpanR pt,0);
    \draw[black!22, rounded corners=4pt]
      (top_y -| \RpanL pt,0) rectangle (bot_y -| \RpanR pt,0);
  \end{scope}

  \end{tikzpicture}

  \caption{A single real Skill, \texttt{d3-visualization}, with the document on
    the left and the code it ships on the right. Each coloured frame encloses
    both ends of one inconsistency; its horizontal span reveals the boundary.
    \textcircled{1}~Both ends sit inside the document (two lines contradict each
    other on the D3 major version), making it an \emph{inter-document}
    inconsistency. \textcircled{2}~The sort rule is stated in the document but
    violated in the shipped code. \textcircled{3}~The deterministic-layout mandate
    appears in the document, but the shipped example uses an unseeded,
    default-tick simulation. \textcircled{2} and \textcircled{3} are
    \emph{document--code} inconsistencies. All three are discoverable by reading
    the files, before any agent runs.}
  \label{fig:motivation}
\end{figure}
